% Generated draft: requires mathematical and source review.
\documentclass[11pt]{article}
\usepackage[margin=25mm]{geometry}
\usepackage{fontspec}
\setmainfont{Latin Modern Roman}
\usepackage{amsmath,amssymb,mathtools,unicode-math}
\setmathfont{Latin Modern Math}
\usepackage{enumitem,needspace,xurl,hyperref}
\setlength{\parindent}{0pt}
\setlength{\parskip}{5pt}
\setlength{\emergencystretch}{3em}
\newcommand{\N}{\mathbb{N}}
\newcommand{\Z}{\mathbb{Z}}
\newcommand{\Q}{\mathbb{Q}}
\newcommand{\R}{\mathbb{R}}
\newcommand{\C}{\mathbb{C}}
\newcommand{\F}{\mathbb{F}}
\newcommand{\A}{\mathbb{A}}
\newcommand{\PP}{\mathbb{P}}
\newcommand{\E}{\mathbb{E}}
\newcommand{\eps}{\varepsilon}
\newcommand{\dd}{\,\mathrm{d}}
\begin{document}
\section*{UnsolvedMath ID: 30006002}
\textbf{Source problem code:} catalogue entry
\subsection{Definitions and mathematical statement}
The supplied record asks for the sharp Bianchi--Egnell stability constant for the classical Sobolev inequality. Its extracted original statement truncates the stability inequality before its right-hand side. The following gives the conventional formulation explicitly; the dimension, scalar field, and distance normalization must be checked against the complete cited report.

Fix an integer \(n\geq 3\), and adopt the real-valued setting. Let \(2^*=2n/(n-2)\), and let \(\dot H^1(\R^n;\R)\) be the completion of \(C_c^\infty(\R^n;\R)\) under the norm
\[
\|u\|_{\dot H^1}=\left(\int_{\R^n}|\nabla u|^2\,\dd x\right)^{1/2},
\]
identified with its representatives in \(L^{2^*}(\R^n)\). Define the sharp Sobolev constant by
\[
S_n=\inf_{\substack{u\in\dot H^1(\R^n;\R)\\u\ne0}}
\frac{\displaystyle\int_{\R^n}|\nabla u|^2\,\dd x}
{\displaystyle\left(\int_{\R^n}|u|^{2^*}\,\dd x\right)^{2/2^*}}.
\]
For \(z\in\R^n\) and \(\lambda>0\), put
\[
B_{z,\lambda}(x)=\lambda^{(n-2)/2}\bigl(1+\lambda^2|x-z|^2\bigr)^{-(n-2)/2},
\qquad
\mathcal M_n=\{aB_{z,\lambda}:a\in\R,\ z\in\R^n,\ \lambda>0\}.
\]
The supplied excerpt identifies constant multiples of these bubbles as the Sobolev optimizers. Set
\[
\delta_n(u)=\int_{\R^n}|\nabla u|^2\,\dd x
-S_n\left(\int_{\R^n}|u|^{2^*}\,\dd x\right)^{2/2^*},
\qquad
d_n(u)^2=\inf_{v\in\mathcal M_n}\int_{\R^n}|\nabla(u-v)|^2\,\dd x.
\]
With this distance convention, define
\[
c_{\mathrm{BE}}(n)=\sup\left\{c\geq0:
\delta_n(u)\geq c\,d_n(u)^2
\text{ for every }u\in\dot H^1(\R^n;\R)\right\}.
\]
Equivalently,
\[
c_{\mathrm{BE}}(n)=
\inf_{\substack{u\in\dot H^1(\R^n;\R)\\d_n(u)>0}}
\frac{\delta_n(u)}{d_n(u)^2}.
\]
The problem is to determine the exact value of \(c_{\mathrm{BE}}(n)\). The source writes \(c_{\mathrm{BE}}\) without selecting a dimension; the dimension-dependent notation here makes that unresolved parameter explicit. The formulation uses all of \(\R^n\), finite homogeneous Sobolev energy, and no finite boundary. The real-valued domain and the displayed distance normalization are explicit reconstruction conventions, since the supplied extraction does not contain the hypotheses and right-hand side of the stability inequality.

The record attributes the existence of a positive stability constant to Bianchi and Egnell. Its supplied dated literature assessment reports improved bounds and no known exact value. These are source-reported claims and status labels, not an independent verification in this draft.
\subsection{Short English statement}
For each dimension \(n\geq3\), determine the largest constant controlling the Sobolev deficit by the squared gradient distance to the family of bubble optimizers. The displayed formulation adopts the conventional real-valued homogeneous Sobolev setting; the incomplete source extraction leaves its exact domain and normalization to be checked. The supplied record labels the problem open.
\subsection{Sources}
\begin{itemize}
\item Local UnsolvedMath registry: \texttt{data/unsolvedmath/problems.json}, permanent ID \(30006002\), title ``Sharp Bianchi--Egnell Stability Constant,'' exact source problem code \texttt{catalogue entry}; supplied collection: Oberwolfach Reports Open Problems. The record gives proposed year \(2024\) and no named proposer.
\item Supplied source citation: \emph{Calculus of Variations (2025)}, Oberwolfach Reports, DOI \texttt{10.4171/owr/2024/37}, pp.\ 2113--2178. \url{https://doi.org/10.4171/owr/2024/37}. The extracted contribution is headed ``The stability of Sobolev's inequality: best constants and optimizers,'' Tobias K\"onig.
\item Supplied literature source: ``Konig, sharp Bianchi--Egnell constant.'' \url{https://londmathsoc.onlinelibrary.wiley.com/doi/abs/10.1112/blms.12837}.
\item Supplied literature source: ``Oberwolfach Report 37/2024.'' \url{https://ems.press/content/serial-article-files/50045}.
\item Attribution retained from the excerpt: Aubin [1] and Talenti [17] for the characterization of Sobolev optimizers; Brezis and Lieb [6] for the motivating stability question; Bianchi and Egnell [3] for quantitative stability. These bracketed labels belong to the source bibliography; their full bibliographic entries were not supplied.
\item Literature assessment checked at \texttt{2026-08-22}: the supplied curation labels the problem open and states that the best constant is positive, improved bounds are available, and no exact value is known for the classical Sobolev inequality. This is supplied dated curation, expressly requiring independent verification of time-sensitive status and cited claims; it has not been independently verified here.
\item The registry labels the statement \texttt{corrected\_verified} and reports reconstruction and checking against the official Oberwolfach report. This is a supplied verification claim. No source was independently visited for this draft.
\end{itemize}
\end{document}
